\subsection{Comparison of Finite Certificates}\label{sec:cert-comparison}
We compare the vertex reduction (V) and the joint Bernstein certificate (B) while keeping
the spline fields, certified levels, and initial covers identical. All variants use
distance-independent refinement with $\vartheta=0.1$ and at most two halvings; screening
uses the unit lift. In addition to $w=1$, we test fixed multipliers with
$\mathbf a=\nabla\phi_A(\mathbf c)$, $\mathbf b=\Rot^\top\nabla\phi_B(\mathbf y_{\mathbf c})$,
\begin{equation}
w_{\rm N}=\sqrt{\frac{\norm{\mathbf b}^2+\delta}{\norm{\mathbf a}^2+\delta}},\qquad
w_{\rm P}=-\frac{\mathbf a^\top\mathbf b}{\norm{\mathbf a}^2+\delta},\quad\delta=10^{-12}.
\end{equation}
The regularizer keeps the weights continuous at zero gradient. Negative projection weights
are permitted, with the absolute value in the remainder of~\eqref{eq:joint-model}.
Neither choice cancels the complete, input-dependent spatial gradient of the CBF condition.

Table~\ref{tab:cert-comparison} uses fixed horizons of 9~s for docking and 6~s for the rounded
tube. The input modification is $J=\int\norm{\mathbf u-\mathbf u^{\rm nom}}\,dt$, compared
only within each scene. All twelve runs require zero slack, and their evaluated distance
lower bounds are positive. The docking runs seat the vehicle; the nominal target is intentionally
10~cm beyond full seating, so reaching that infeasible target is not the docking success criterion.

\begin{table}[t]
\caption{Certificate Comparison with Fixed Geometry}
\label{tab:cert-comparison}
\centering\footnotesize
\setlength{\tabcolsep}{3pt}
\begin{tabular}{lrrrr}
\toprule
Certificate & Dock $J$ & Tube $t_g$ [s] & Tube $J$ & Tube time [ms]\\
\midrule
V, $w=1$       & 0.656 & 3.99 & 1.948 & 34.6/83.7\\
B, $w=1$       & 0.653 & 3.93 & 1.870 & 52.2/165.7\\
V, $w_{\rm N}$ & 0.655 & 3.99 & 1.951 & 33.9/85.3\\
B, $w_{\rm N}$ & 0.653 & 3.92 & 1.857 & 52.2/165.6\\
V, $w_{\rm P}$ & 0.655 & 4.00 & 1.965 & 34.7/83.3\\
B, $w_{\rm P}$ & 0.653 & 3.95 & 1.896 & 51.0/166.5\\
\bottomrule
\end{tabular}
\\[3pt]
\parbox{\columnwidth}{\footnotesize $t_g$: simulated goal-arrival time. Filter time: median/p95
over the fixed 6-s horizon, one CPU thread. These are exploratory measurements with background
activity not excluded, rather than the isolated timings of Table~\ref{tab:franka}. The tube's
15-s experiment above has a different timing distribution.}
\end{table}

The joint unit-weight certificate reduces tube input modification by about 4\%, at the cost
of increasing the maximum row count from 26{,}440 to 89{,}424. The fixed multipliers do not
improve the vertex variant's tube arrival time or input modification. These cases therefore
support retaining the unit-weight vertex reduction as the economical implementation, while
the joint coefficient condition provides a direct alternative when the additional cost is
acceptable. The multiplier tests establish neither a general improvement nor a second-order
conservatism result.

We also compare the two unit-weight certificates on six saved Franka trials, selected as three
that reached and three that stalled in the original experiment. With the same distance-independent
refinement, both variants reach in four of six trials and require zero slack in all six.
For the newly successful trial, the joint certificate reduces arrival time from 7.38 to 6.63~s
and $J$ from 4.243 to 3.533; its median/p95 filter time increases from 46/116 to 73/250~ms.
In the most expensive trial neither variant reaches within 10~s, while median filter time rises
from 115 to 253~ms and maximum row count from 41{,}776 to 146{,}691. These selected trials
show a cost--performance tradeoff, not a new estimate of the 30-trial success rate.
